\documentclass[11pt]{article}

\usepackage[T1]{fontenc}
\usepackage[utf8]{inputenc}
\usepackage{lmodern}
\usepackage[margin=1.15in]{geometry}
\usepackage{amsmath,amssymb,amsthm,mathtools}
\usepackage{enumitem}
\usepackage{booktabs}
\usepackage{hyperref}
\usepackage{xcolor}

\hypersetup{
  colorlinks=true,
  linkcolor=blue!50!black,
  citecolor=blue!50!black,
  urlcolor=blue!50!black
}

\newtheorem{theorem}{Theorem}[section]
\newtheorem{lemma}[theorem]{Lemma}
\newtheorem{corollary}[theorem]{Corollary}
\newtheorem{proposition}[theorem]{Proposition}
\newtheorem{definition}[theorem]{Definition}
\newtheorem{remark}[theorem]{Remark}

\newcommand{\N}{\mathbb{N}}
\newcommand{\ceil}[1]{\left\lceil #1 \right\rceil}
\newcommand{\floor}[1]{\left\lfloor #1 \right\rfloor}
\newcommand{\gdiamond}{g^{\diamond}}
\newcommand{\astar}{g^*}
\newcommand{\alphaQ}{\alpha_Q}
\newcommand{\alphaJQ}{\alpha_J^Q}
\newcommand{\alphaJS}{\alpha_J^S}

\title{A Direct Ackermann Comparison for the Seidel--Sharir Top-Down Path-Compression Hierarchy}
\author{Path-compression proof-digestion note}
\date{May 2026}

\begin{document}
\maketitle

\begin{abstract}
This note gives a source-faithful, proof-digested presentation of how the Seidel--Sharir top-down path-compression hierarchy yields an inverse-Ackermann bound under an explicit Ackermann normalization.  The central comparison is phrased through the integer-threshold inverse
\[
R_k(t)=\max\{r\ge 0:J_k(r)\le t\}\qquad(t\in\N),
\]
and proves
\[
R_{z+1}(Q)\ge A(z,4Q)
\]
for all integers $z\ge 1$ and $Q\ge 1$.  Thus the comparison target $R_{cz+C}(Q)\ge A(z,DQ)$ holds with constants $c=1$, $C=1$, and $D=4$.  The proof stays within the Seidel--Sharir top-down $J$-hierarchy and the associated $g^\diamond$ recurrence; it does not use a bottom-up Tarjan-potential argument.  Translating the comparison to the packet normalizations $L(n)=\ceil{\log_2\max(n,2)}$ and $Q(m,n)=\ceil{1+m/n}$ gives $\alphaJQ(m,n)\le \alphaQ(m,n)+1$, the source-faithful real-threshold bound $\alphaJS(m,n)\le\alphaQ(m,n)+2$ with a $+1$ improvement in the integral case, and the cost bound
\[
f(m,n,L(n))\le (\alphaQ(m,n)+3)m+4n.
\]
\end{abstract}

\section{Introduction}

Seidel and Sharir gave a top-down analysis of path compression in union-find structures, deriving recurrences from which the usual bounds arise naturally; their abstract emphasizes that the inverse-Ackermann bound can be derived without introducing the Ackermann function itself \cite{SeidelSharir2005}.  The Dagstuhl Seminar 25191 report records Tarjan's later comparison question about the top-down analysis in Section~5.3 of the seminar report \cite{BenderEtAlDagstuhl25191}.  Classical union-find analysis and the Ackermann hierarchy enter through Tarjan's work on set union \cite{Tarjan1975}.

This note addresses that comparison problem at the level of proof exposition and digestion.  It gives a direct comparison from the Seidel--Sharir top-down $J$-hierarchy to a concrete Ackermann normalization.  The point is not to claim a new union-find upper bound.  Rather, the contribution is a source-anchored proof that the Seidel--Sharir $J$ hierarchy itself implies an inverse-Ackermann bound, with all constants and normalizations explicit.

The main theorem is the threshold comparison
\[
R_{z+1}(Q)\ge A(z,4Q)
\qquad (z\ge 1,\ Q\ge 1),
\]
where, for integer thresholds $t\ge0$,
\[
R_k(t)=\max\{r\ge 0:J_k(r)\le t\}.
\]
Thus the target comparison $R_{cz+C}(Q)\ge A(z,DQ)$ holds with
\[
c=1,\qquad C=1,\qquad D=4.
\]
The rest of the paper translates this comparison into the packet-normalized alpha bounds and the resulting cost estimate.

\paragraph{What is new here?}
This note does not claim a new union-find bound.  It gives a direct, source-anchored proof that the Seidel--Sharir top-down $J$ hierarchy implies an inverse-Ackermann bound under an explicit Ackermann normalization, addressing Tarjan's Dagstuhl comparison problem at the level of proof exposition and digestion.

\section{Source setup and conventions}

This section records the source facts and packet normalizations used in the proof.  The source facts are the top-down path-compression recurrence machinery of Seidel and Sharir \cite{SeidelSharir2005}; the later proof only compares the resulting $J$ hierarchy with an Ackermann hierarchy.

\subsection{The transforms \texorpdfstring{$g^*$ and $g^\diamond$}{g* and g-diamond}}

For an integer-valued function $g$ satisfying $g(r)<r$ for $r>0$, Seidel and Sharir define an iteration count $g^*$ by
\[
g^*(r)=
\begin{cases}
0, & r\le 1,\\
1+g^*(g(r)), & r>1,
\end{cases}
\]
and use the related transform
\[
g^\diamond(r)=
\begin{cases}
g(r), & g(r)\le 1,\\[4pt]
1+g^\diamond(\ceil{\log_2 g(r)}), & g(r)>1.
\end{cases}
\]
All logarithms in this paper are base $2$.

\subsection{The \texorpdfstring{$J$}{J} hierarchy}

Seidel--Sharir Corollary~6 defines the hierarchy
\[
J_0(r)=\ceil{\frac{r-1}{2}},
\qquad
J_{k+1}=J_k^\diamond.
\]
Equivalently,
\[
J_{k+1}(r)=
\begin{cases}
J_k(r),&J_k(r)\le 1,\\[4pt]
1+J_{k+1}(\ceil{\log_2 J_k(r)}),&J_k(r)>1.
\end{cases}
\]

\subsection{The source-correct shifting lemma and recurrence}

The source-correct top-down shifting lemma is Seidel--Sharir Lemma~5:
\begin{quote}
If
\[
f(m,n,r)\le km+2n g(r),
\]
then
\[
f(m,n,r)\le (k+1)m+2n g^\diamond(r).
\]
\end{quote}
Iterating this lemma through the $J$ hierarchy gives the source recurrence consequence
\[
f(m,n,r)\le km+2nJ_k(r).
\]
The factor $2n$ is part of the source statement and is used throughout the cost translation.

\subsection{Rank and threshold normalizations}

The source rank cutoff is governed by the rank bound $\log_2 n$.  The packet convention stabilizes the cutoff as
\[
L(n)=\ceil{\log_2\max(n,2)}.
\]
For $n\ge 2$, this agrees with $\ceil{\log_2 n}$; for $n=1$, it avoids a degenerate zero cutoff.

The source inverse threshold is
\[
J_k(\ceil{\log_2 n})\le 1+\frac mn.
\]
The packet integer threshold is
\[
Q(m,n)=\ceil{1+\frac mn}.
\]
This is a normalization device: it turns the real threshold $1+m/n$ into an integer threshold suitable for the threshold inverse $R_k$.

\begin{remark}[The non-source-faithful $2m/n$ threshold]
The threshold $J_k(\lg n)\le 2m/n$ is not the source threshold recorded in the current packet.  It is treated below only as a quarantined alternate target; the source-faithful threshold is $1+m/n$.
\end{remark}

\section{Blackboard proof idea}

The proof has one moving part: the $\diamond$ recurrence turns a bound at level $k$ into a much larger threshold inverse at level $k+1$.

For integer thresholds $t\ge0$, define
\[
R_k(t)=\max\{r\ge 0:J_k(r)\le t\}.
\]
If $r\le R_k(2^B)$, then $J_k(r)\le 2^B$.  In the recursive branch of $J_{k+1}=J_k^\diamond$,
\[
J_{k+1}(r)=1+J_{k+1}(\ceil{\log_2J_k(r)}),
\]
and the logarithmic argument is at most $B$.  Therefore, if $B\le R_{k+1}(t)$, then the recursive call costs at most $t$, and
\[
J_{k+1}(r)\le t+1.
\]
This gives the threshold recurrence
\[
R_{k+1}(t+1)\ge R_k(2^{R_{k+1}(t)}).
\]
That recurrence has Ackermann flavor: increasing the $J$-level by one lets us feed an exponential of a previous threshold into the lower-level inverse.

The proof turns this into the invariant
\[
R_j(x)\ge A(j+1,x)
\qquad(j\ge1,\ x\ge1).
\]
The base level $R_0(t)=2t+1$ is too small to dominate $A(1,t)=2^t$; this is why the comparison has a one-level shift.  After the invariant is proved, a threshold jump gives
\[
R_{z+1}(Q)\ge A(z,4Q).
\]

\section{Technical proof}

\subsection{Ackermann normalization and threshold inverse}

\begin{definition}[Ackermann normalization]
Define $A(i,x)$ by
\[
A(0,x)=2x,
\]
and, for $i\ge 1$,
\[
A(i,0)=0,
\qquad
A(i,1)=2,
\qquad
A(i,x)=A(i-1,A(i,x-1))\quad(x\ge2).
\]
In particular, $A(1,x)=2^x$ for $x\ge1$.
\end{definition}

\begin{definition}[Threshold inverse]
For an integer threshold $t\ge0$, define
\[
R_k(t)=\max\{r\ge0:J_k(r)\le t\}.
\]
The maximum is finite by Lemma~\ref{lem:J-package} below.
\end{definition}

\subsection{The basic \texorpdfstring{$J_k$}{Jk} package}

\begin{lemma}[Diamond preservation]\label{lem:diamond-preservation}
Let $g:\N\to\N$ be integer-valued, nonnegative, nondecreasing, unbounded, and satisfy
\[
g(0)=0,
\qquad
g(r)<r\quad(r>0).
\]
Define $h=g^\diamond$ by
\[
h(r)=
\begin{cases}
g(r),&g(r)\le1,\\[4pt]
1+h(\ceil{\log_2 g(r)}),&g(r)>1.
\end{cases}
\]
Then $h$ is integer-valued, nonnegative, nondecreasing, unbounded, satisfies $h(0)=0$, and satisfies
\[
h(r)\le g(r)\quad(r\ge0),
\]
hence
\[
h(r)<r\quad(r>0).
\]
\end{lemma}

\begin{proof}
For every integer $x\ge2$,
\[
\ceil{\log_2 x}\le x-1.
\]
Thus, in the recursive branch $g(r)>1$, with $u=\ceil{\log_2 g(r)}$, we have
\[
u\le g(r)-1<r.
\]
So $h(r)$ is defined by recursion on smaller arguments.  It is integer-valued and nonnegative because $g$ is.  Also $g(1)=0$, since $g(1)<1$ and $g$ is nonnegative.

We prove $h(r)\le g(r)$ by induction on $r$.  If $g(r)\le1$, then $h(r)=g(r)$.  If $g(r)>1$, set $u=\ceil{\log_2 g(r)}$.  Since $u<r$, induction gives $h(u)\le g(u)$.  Since $g(0)=0$ and $g(u)<u$ for $u>0$, we have $g(u)\le u$ for every $u$.  Thus
\[
h(r)=1+h(u)\le1+u\le g(r).
\]
This proves $h(r)\le g(r)$, and hence $h(r)<r$ for $r>0$.

It remains to prove monotonicity.  We prove by strong induction on $s$ that $r\le s$ implies $h(r)\le h(s)$.  Let $r\le s$, and set $x=g(r)$ and $y=g(s)$.  Since $g$ is nondecreasing, $x\le y$.  If $y\le1$, then $h(r)=x\le y=h(s)$.  If $y>1$ but $x\le1$, then $h(r)=x\le1\le h(s)$.  If $x>1$, put
\[
u=\ceil{\log_2 x},
\qquad
v=\ceil{\log_2 y}.
\]
Then $u\le v$ and $v\le y-1<s$.  By strong induction, $h(u)\le h(v)$.  Therefore
\[
h(r)=1+h(u)\le1+h(v)=h(s).
\]
Thus $h$ is nondecreasing.

Finally, $h$ is unbounded.  We prove by induction on $M\ge0$ that there exists $r$ with $h(r)\ge M$.  For $M=0$ this is immediate.  Suppose it holds for $M$, and choose $u$ with $h(u)\ge M$.  Since $g$ is unbounded, choose $r$ with $g(r)>2^u$.  Then $g(r)>1$ and $v=\ceil{\log_2 g(r)}>u$.  By monotonicity, $h(v)\ge h(u)\ge M$.  Hence
\[
h(r)=1+h(v)\ge M+1.
\]
This proves unboundedness.
\end{proof}

\begin{lemma}[Basic $J_k$ package]\label{lem:J-package}
For every $k\ge0$:
\begin{enumerate}[label=(\roman*)]
\item $J_k$ is integer-valued and nonnegative on $\N$;
\item $J_k(0)=0$;
\item $J_k$ is nondecreasing;
\item $J_k$ is unbounded on $\N$;
\item $J_k(r)<r$ for every $r>0$, in particular for every $r>1$;
\item $J_{k+1}(r)\le J_k(r)$ for every $r\ge0$;
\item $R_k(t)$ is finite for every integer $t\ge0$;
\item for integer thresholds $0\le t\le t'$, we have $R_k(t)\le R_k(t')$, and $R_{k+1}(t)\ge R_k(t)$.
\end{enumerate}
\end{lemma}

\begin{proof}
For $k=0$,
\[
J_0(r)=\ceil{\frac{r-1}{2}}.
\]
This is integer-valued and nonnegative for integer $r\ge0$.  It satisfies $J_0(0)=0$, is nondecreasing, is unbounded, and satisfies $J_0(r)<r$ for $r>0$.

Assume $J_k$ has the first five properties.  Apply Lemma~\ref{lem:diamond-preservation} to $g=J_k$.  Since $J_{k+1}=J_k^\diamond$, it follows that $J_{k+1}$ is integer-valued, nonnegative, nondecreasing, unbounded, satisfies $J_{k+1}(0)=0$, and satisfies
\[
J_{k+1}(r)\le J_k(r)<r\quad(r>0).
\]
This proves the first six claims by induction on $k$.

Fix an integer threshold $t\ge0$.  Since $J_k$ is unbounded and nondecreasing, the set $\{r\ge0:J_k(r)\le t\}$ is finite.  Since $J_k(0)=0\le t$, it is nonempty.  Hence $R_k(t)$ is a finite maximum.  If $0\le t\le t'$ are integers, then $\{r:J_k(r)\le t\}\subseteq\{r:J_k(r)\le t'\}$, so $R_k(t)\le R_k(t')$.  Since $J_{k+1}(r)\le J_k(r)$, we also have $R_{k+1}(t)\ge R_k(t)$.
\end{proof}

\subsection{Threshold recurrences}

\begin{lemma}[Exact base inverse]\label{lem:R0}
For every integer $t\ge0$,
\[
R_0(t)=2t+1.
\]
\end{lemma}

\begin{proof}
For integer $t\ge0$, the inequality $J_0(r)\le t$ is equivalent to
\[
\ceil{\frac{r-1}{2}}\le t,
\]
which is equivalent to $r\le2t+1$.  Hence the largest such $r$ is $2t+1$.
\end{proof}

\begin{lemma}[Diamond-to-threshold recurrence]\label{lem:diamond-threshold}
Let $B,t$ be nonnegative integers.  If
\[
r\le R_k(2^B)
\]
and
\[
B\le R_{k+1}(t),
\]
then
\[
J_{k+1}(r)\le t+1.
\]
Consequently,
\[
R_{k+1}(t+1)\ge R_k(2^{R_{k+1}(t)}).
\]
\end{lemma}

\begin{proof}
Assume $r\le R_k(2^B)$.  Then $J_k(r)\le2^B$.  If $J_k(r)\le1$, then $J_{k+1}(r)=J_k(r)\le1\le t+1$.  If $J_k(r)>1$, then
\[
J_{k+1}(r)=1+J_{k+1}(\ceil{\log_2J_k(r)}).
\]
Since $J_k(r)\le2^B$, we have $\ceil{\log_2J_k(r)}\le B$.  By Lemma~\ref{lem:J-package}, $J_{k+1}$ is nondecreasing.  Since $B\le R_{k+1}(t)$,
\[
J_{k+1}(\ceil{\log_2J_k(r)})\le J_{k+1}(B)\le t.
\]
Thus $J_{k+1}(r)\le t+1$.

Therefore every $r\le R_k(2^B)$ satisfies $J_{k+1}(r)\le t+1$, so $R_{k+1}(t+1)\ge R_k(2^B)$.  Taking $B=R_{k+1}(t)$ proves the displayed recurrence.
\end{proof}

\begin{lemma}[Threshold jump]\label{lem:threshold-jump}
For every $k\ge0$ and every $Q\ge2$,
\[
R_{k+1}(Q)\ge R_k(4Q).
\]
\end{lemma}

\begin{proof}
Apply Lemma~\ref{lem:diamond-threshold} with $t=Q-1$:
\[
R_{k+1}(Q)\ge R_k(2^{R_{k+1}(Q-1)}).
\]
By Lemmas~\ref{lem:J-package} and~\ref{lem:R0},
\[
R_{k+1}(Q-1)\ge R_0(Q-1)=2Q-1.
\]
For $Q\ge2$,
\[
2^{2Q-1}\ge4Q.
\]
Using monotonicity of $R_k$ in its threshold argument,
\[
R_{k+1}(Q)\ge R_k(2^{R_{k+1}(Q-1)})\ge R_k(4Q).
\]
\end{proof}

\subsection{Ackermann facts}

\begin{lemma}[Ackermann monotonicity and row domination]\label{lem:ackermann-basic}
For the Ackermann function $A$ defined above:
\begin{enumerate}[label=(\roman*)]
\item for every $i\ge0$, $A(i,x)$ is nondecreasing in $x$;
\item for every $i\ge1$ and $x\ge1$, $A(i,x)\ge2x$;
\item for every $i\ge0$ and $x\ge1$, $A(i+1,x)\ge A(i,x)$.
\end{enumerate}
\end{lemma}

\begin{proof}
We first prove monotonicity in the second argument and the growth bound together by induction on the row index.  For $i=0$, monotonicity follows from $A(0,x)=2x$.

Fix $i\ge1$ and assume the corresponding lower-row facts have already been proved.  We prove row-$i$ monotonicity by induction on $x$.  Since $A(i,0)=0$ and $A(i,1)=2$, the first step is clear.  For $x\ge1$,
\[
A(i,x+1)=A(i-1,A(i,x)).
\]
If $i=1$, then $A(0,y)=2y\ge y$ for $y\ge0$.  If $i\ge2$, the lower-row growth bound gives $A(i-1,y)\ge 2y\ge y$ for $y\ge1$.  Hence $A(i,x+1)\ge A(i,x)$, and row $i$ is nondecreasing.

We next prove the growth bound for row $i$.  At $x=1$, $A(i,1)=2=2x$.  For $x\ge2$, again using
\[
A(i,x)=A(i-1,A(i,x-1)),
\]
we get $A(i,x)=2A(i,x-1)$ when $i=1$, and $A(i,x)\ge2A(i,x-1)$ when $i\ge2$ by the lower-row growth bound.  Thus, by induction on $x$ with base $A(i,1)=2$,
\[
A(i,x)\ge2A(i,x-1)\ge4(x-1)\ge2x.
\]
This proves monotonicity and growth for all rows.

For row domination, fix $i\ge0$.  At $x=1$, $A(i+1,1)=2=A(i,1)$.  For $x\ge2$,
\[
A(i+1,x)=A(i,A(i+1,x-1)).
\]
The growth bound gives $A(i+1,x-1)\ge2(x-1)\ge x$, and row $i$ is nondecreasing.  Hence $A(i+1,x)\ge A(i,x)$.
\end{proof}

\begin{lemma}[Ackermann domination invariant]\label{lem:ackermann-domination}
For every $j\ge1$ and $x\ge1$,
\[
R_j(x)\ge A(j+1,x).
\]
\end{lemma}

\begin{proof}
We prove the claim by induction on $j$, with an inner induction on $x$.

First let $j=1$.  For $x=1$, Lemma~\ref{lem:diamond-threshold} with $k=0,t=0$ gives
\[
R_1(1)\ge R_0(2^{R_1(0)}).
\]
By Lemma~\ref{lem:J-package}, $R_1(0)\ge R_0(0)=1$.  Thus
\[
R_1(1)\ge R_0(2)=5\ge2=A(2,1).
\]
Assume $R_1(x)\ge A(2,x)$.  By Lemma~\ref{lem:diamond-threshold},
\[
R_1(x+1)\ge R_0(2^{R_1(x)}).
\]
Since $R_0(t)=2t+1$,
\[
R_0(2^{R_1(x)})=2\cdot2^{R_1(x)}+1.
\]
Using the induction hypothesis and $A(1,y)=2^y$,
\[
2\cdot2^{R_1(x)}+1\ge2^{A(2,x)}=A(2,x+1).
\]
Thus the case $j=1$ holds.

Now assume $R_j(x)\ge A(j+1,x)$ for all $x\ge1$.  We prove $R_{j+1}(x)\ge A(j+2,x)$ for all $x\ge1$.  For $x=1$, Lemma~\ref{lem:J-package} gives
\[
R_{j+1}(1)\ge R_0(1)=3\ge2=A(j+2,1).
\]
Assume $R_{j+1}(x)\ge A(j+2,x)$.  By Lemma~\ref{lem:diamond-threshold},
\[
R_{j+1}(x+1)\ge R_j(2^{R_{j+1}(x)}).
\]
By the outer induction hypothesis,
\[
R_j(y)\ge A(j+1,y)\qquad(y\ge1).
\]
Thus
\[
R_j(2^{R_{j+1}(x)})\ge A(j+1,2^{R_{j+1}(x)}).
\]
By the inner inductive hypothesis,
\[
R_{j+1}(x)\ge A(j+2,x).
\]
Since $2^k>k$ for every integer $k\ge0$, we have
\[
2^{R_{j+1}(x)}>R_{j+1}(x)\ge A(j+2,x).
\]
By Lemma~\ref{lem:ackermann-basic}(i), $A(j+1,-)$ is nondecreasing, so
\[
A(j+1,2^{R_{j+1}(x)})\ge A(j+1,A(j+2,x))=A(j+2,x+1),
\]
where the last equality is the Ackermann recursion.  Hence $R_{j+1}(x+1)\ge A(j+2,x+1)$.
\end{proof}

\subsection{Main comparison theorem}

\begin{theorem}[Main comparison]\label{thm:main-comparison}
For all integers $z\ge1$ and $Q\ge1$,
\[
R_{z+1}(Q)\ge A(z,4Q).
\]
Consequently, for every integer $r\ge0$,
\[
A(z,4Q)>r\quad\Longrightarrow\quad J_{z+1}(r)\le Q.
\]
\end{theorem}

\begin{proof}
First suppose $Q\ge2$.  By Lemma~\ref{lem:threshold-jump},
\[
R_{z+1}(Q)\ge R_z(4Q).
\]
Since $z\ge1$, Lemma~\ref{lem:ackermann-domination} gives
\[
R_z(4Q)\ge A(z+1,4Q).
\]
By row domination from Lemma~\ref{lem:ackermann-basic},
\[
A(z+1,4Q)\ge A(z,4Q).
\]
Therefore $R_{z+1}(Q)\ge A(z,4Q)$.

Now suppose $Q=1$ and $z=1$.  Lemma~\ref{lem:diamond-threshold} gives
\[
R_2(1)\ge R_1(2^{R_2(0)}).
\]
By Lemma~\ref{lem:J-package}, $R_2(0)\ge R_0(0)=1$, so $R_2(1)\ge R_1(2)$.  The proof of Lemma~\ref{lem:ackermann-domination} gave $R_1(1)\ge5$, and Lemma~\ref{lem:diamond-threshold} gives
\[
R_1(2)\ge R_0(2^{R_1(1)})\ge R_0(2^5)=65.
\]
Thus
\[
R_2(1)\ge65\ge16=A(1,4).
\]

Finally suppose $Q=1$ and $z\ge2$.  Lemma~\ref{lem:diamond-threshold} gives
\[
R_{z+1}(1)\ge R_z(2^{R_{z+1}(0)}).
\]
By Lemma~\ref{lem:J-package}, $R_{z+1}(0)\ge R_0(0)=1$, so $R_{z+1}(1)\ge R_z(2)$.  Apply Lemma~\ref{lem:threshold-jump} with level $z-1$ and $Q=2$:
\[
R_z(2)\ge R_{z-1}(8).
\]
Because $z-1\ge1$, Lemma~\ref{lem:ackermann-domination} gives
\[
R_{z-1}(8)\ge A(z,8).
\]
By Ackermann monotonicity,
\[
A(z,8)\ge A(z,4).
\]
Therefore $R_{z+1}(1)\ge A(z,4)$.

The three cases prove $R_{z+1}(Q)\ge A(z,4Q)$ for all $z\ge1$ and $Q\ge1$.  If $A(z,4Q)>r$, then $r\le A(z,4Q)-1\le R_{z+1}(Q)$, and therefore $J_{z+1}(r)\le Q$.
\end{proof}

\section{Alpha consequences and cost bound}

\begin{definition}[Packet inverses]
Define
\[
\alphaJQ(m,n)=\min\{k\ge0:J_k(L(n))\le Q(m,n)\},
\]
and
\[
\alphaQ(m,n)=\min\{z\ge1:A(z,4Q(m,n))>L(n)\}.
\]
\end{definition}

\begin{lemma}[Existence of $\alphaQ$]\label{lem:alphaQ-exists}
For all positive integers $m,n$, the minimum defining $\alphaQ(m,n)$ exists.
\end{lemma}

\begin{proof}
Since $Q(m,n)\ge1$, we have $4Q(m,n)\ge4$.  Let $a_z=A(z,4)$.  Then $a_1=16$, and for $z\ge1$,
\[
A(z+1,2)=4,
\qquad
A(z+1,3)=A(z,4)=a_z,
\]
so
\[
a_{z+1}=A(z+1,4)=A(z,a_z).
\]
By the growth bound in Lemma~\ref{lem:ackermann-basic}, $A(z,a_z)\ge2a_z$.  Hence $a_z$ is unbounded in $z$.  By monotonicity in the second argument,
\[
A(z,4Q(m,n))\ge A(z,4)=a_z,
\]
so some $z\ge1$ satisfies $A(z,4Q(m,n))>L(n)$.
\end{proof}

\begin{theorem}[Patched threshold consequence]\label{thm:patched-alpha}
For all positive integers $m,n$,
\[
\alphaJQ(m,n)\le\alphaQ(m,n)+1.
\]
\end{theorem}

\begin{proof}
Let $z=\alphaQ(m,n)$, $Q=Q(m,n)$, and $L=L(n)$.  By definition,
\[
A(z,4Q)>L.
\]
By Theorem~\ref{thm:main-comparison},
\[
R_{z+1}(Q)\ge A(z,4Q)>L.
\]
Since $L$ is an integer, $L\le R_{z+1}(Q)$.  Therefore $J_{z+1}(L)\le Q$, and hence
\[
\alphaJQ(m,n)\le z+1=\alphaQ(m,n)+1.
\]
\end{proof}

\begin{definition}[Source-faithful real-threshold inverse]
Define
\[
\alphaJS(m,n)=\min\left\{k\ge0:J_k(L(n))\le1+\frac mn\right\}.
\]
\end{definition}

\begin{lemma}[Ackermann buffer]\label{lem:ackermann-buffer}
For $z\ge1$ and integer $p\ge1$,
\[
A(z+1,4p)\ge A(z,4p+4).
\]
\end{lemma}

\begin{proof}
Since $4p\ge2$,
\[
A(z+1,4p)=A(z,A(z+1,4p-1)).
\]
We claim $A(z+1,4p-1)\ge4p+4$.  If $p=1$, then $4p-1=3$, and $A(i,2)=4$ for all $i\ge1$, so
\[
A(z+1,3)=A(z,4)\ge8=4p+4.
\]
If $p\ge2$, Lemma~\ref{lem:ackermann-basic} gives
\[
A(z+1,4p-1)\ge2(4p-1)=8p-2\ge4p+4.
\]
By monotonicity of $A(z,-)$, $A(z+1,4p)\ge A(z,4p+4)$.
\end{proof}

\begin{theorem}[Source-faithful real-threshold consequence]\label{thm:source-alpha}
For all positive integers $m,n$,
\[
\alphaJS(m,n)\le\alphaQ(m,n)+2.
\]
If $1+m/n$ is integral, then
\[
\alphaJS(m,n)\le\alphaQ(m,n)+1.
\]
\end{theorem}

\begin{proof}
Let
\[
s=1+\frac mn,
\qquad
Q=\ceil{s},
\qquad
L=L(n),
\qquad
z=\alphaQ(m,n).
\]
Then $A(z,4Q)>L$.

If $s$ is an integer, then $Q=s$.  Theorem~\ref{thm:main-comparison} gives
\[
R_{z+1}(s)=R_{z+1}(Q)\ge A(z,4Q)>L.
\]
Thus $J_{z+1}(L)\le s$, and $\alphaJS(m,n)\le z+1=\alphaQ(m,n)+1$.

If $s$ is not integral, put $p=\floor{s}$.  Since $m,n\ge1$, $p\ge1$, and $Q=p+1$.  Thus $A(z,4p+4)>L$.  By Lemma~\ref{lem:ackermann-buffer}, $A(z+1,4p)>L$.  Theorem~\ref{thm:main-comparison} gives
\[
R_{z+2}(p)\ge A(z+1,4p)>L.
\]
Thus $J_{z+2}(L)\le p<s$, and therefore
\[
\alphaJS(m,n)\le z+2=\alphaQ(m,n)+2.
\]
\end{proof}

\begin{theorem}[Cost consequence]\label{thm:cost}
Using the source recurrence $f(m,n,r)\le km+2nJ_k(r)$, one obtains
\[
f(m,n,L(n))\le(\alphaQ(m,n)+3)m+4n.
\]
Consequently,
\[
f(m,n,L(n))=O(m\alphaQ(m,n)+n).
\]
\end{theorem}

\begin{proof}
Let $z=\alphaQ(m,n)$, $Q=Q(m,n)$, and $L=L(n)$.  By Theorem~\ref{thm:patched-alpha}, $J_{z+1}(L)\le Q$.  Applying the source recurrence with $k=z+1$ and $r=L$ gives
\[
f(m,n,L)\le(z+1)m+2nJ_{z+1}(L)\le(z+1)m+2nQ.
\]
Since
\[
Q=\ceil{1+\frac mn}\le2+\frac mn,
\]
we have
\[
2nQ\le2n\left(2+\frac mn\right)=4n+2m.
\]
Therefore
\[
f(m,n,L)\le(z+3)m+4n.
\]
Substituting $z=\alphaQ(m,n)$ proves the claim.
\end{proof}

\section{Normalization discussion}

\begin{definition}[Classical row inverse for this normalization]
Define
\[
\alpha_A(N,X)=\min\{z\ge1:A(z,X)>N\}.
\]
\end{definition}
Then the packet inverse is exactly
\[
\alphaQ(m,n)=\alpha_A(L(n),4Q(m,n)).
\]
The proved comparison says
\[
\alphaJQ(m,n)\le \alpha_A(L(n),4Q(m,n))+1.
\]
For the source-faithful real threshold,
\[
\alphaJS(m,n)\le \alpha_A(L(n),4Q(m,n))+2,
\]
with the $+1$ improvement when $1+m/n$ is integral.

The source defines
\[
\alpha_S(m,n)=\min\left\{k:J_k(\ceil{\log_2 n})\le1+\frac mn\right\}.
\]
For $n\ge2$, $L(n)=\ceil{\log_2 n}$, so $\alphaJS$ is the packet notation for this source threshold.  For $n=1$, the packet uses $L(1)=1$ rather than $0$ as a harmless stabilization.

The Seidel--Sharir paper discusses Tarjan's initial inverse-Ackermann function $\alpha_T$, defined through a different hierarchy $T_k$ \cite{SeidelSharir2005,Tarjan1975}.  This paper does not use that comparison as a proof dependency.  The theorem here is restricted to the concrete normalization $\alphaQ(m,n)=\alpha_A(L(n),4Q(m,n))$.

\section{Hazards avoided}

\subsection{The threshold \texorpdfstring{$1+m/n$ versus $2m/n$}{1+m/n versus 2m/n}}

The source-faithful $J$ threshold is
\[
J_k(\ceil{\log_2 n})\le1+\frac mn,
\]
not a $J_k$ threshold $2m/n$.  The alternate condition
\[
J_k(\lg n)\le\frac{2m}{n}
\]
is not source-faithful under the current packet unless independently verified.

Moreover, it is impossible for $m/n<1/2$ and $n\ge3$.  Indeed, $n\ge3$ gives $L(n)\ge2$, and a direct induction on $k$ gives $J_k(r)\ge1$ for every $r\ge2$: the claim is clear for $J_0$, and the $g^\diamond$ definition either returns $1$ in the small branch or $1+$ something in the recursive branch.  If $m/n<1/2$, then $2m/n<1$, so the alternate threshold would force $J_k(L(n))=0$, contradicting $J_k(L(n))\ge1$.

\subsection{The factor \texorpdfstring{$2n$}{2n} in the shifting lemma}

The source shifting lemma has premise
\[
f(m,n,r)\le km+2ng(r)
\]
and conclusion
\[
f(m,n,r)\le(k+1)m+2n g^\diamond(r).
\]
The earlier shorthand with $ng(r)$ is not used here.

\subsection{Why \texorpdfstring{$Q(m,n)=\ceil{1+m/n}$}{Q(m,n)=ceil(1+m/n)} is used}

The comparison theorem is stated for an integer threshold $Q\ge1$, while the source threshold $1+m/n$ is real-valued.  The packet threshold
\[
Q(m,n)=\ceil{1+\frac mn}
\]
provides an integer target for the $R_k$ comparison.  The exact real source threshold is recovered with one extra buffer in the nonintegral case.

\subsection{Avoiding bottom-up Tarjan-potential drift}

The proof uses only the Seidel--Sharir top-down $J$ hierarchy, the $g^\diamond$ threshold recurrence, the packet Ackermann normalization, and the source recurrence $f(m,n,r)\le km+2nJ_k(r)$.  It does not use bottom-up Tarjan potentials.

\section{Conclusion}

The direct bridge is
\[
R_{z+1}(Q)\ge A(z,4Q),
\]
with constants $c=1$, $C=1$, and $D=4$ in the comparison target.  This bridge converts the Seidel--Sharir $J$ hierarchy into the packet-normalized inverse-Ackermann bound.  The result is proof-digestion and exposition: it does not claim a new union-find upper bound, but it makes the Ackermann comparison explicit and source-faithful.  It is presented as a source-faithful answer to Tarjan's Dagstuhl 25191 \S5.3 comparison prompt rather than a priority claim.

\begin{corollary}[Final citable theorem block]\label{cor:final-block}
Assume the Seidel--Sharir $J$ hierarchy and the packet Ackermann normalization above.  Then, for all integers $z\ge1$ and $Q\ge1$,
\[
R_{z+1}(Q)\ge A(z,4Q).
\]
Thus the comparison target $R_{cz+C}(Q)\ge A(z,DQ)$ holds with
\[
c=1,\qquad C=1,\qquad D=4.
\]
Let
\[
L(n)=\ceil{\log_2\max(n,2)},
\qquad
Q(m,n)=\ceil{1+\frac mn},
\]
and define $\alphaQ$, $\alphaJQ$, and $\alphaJS$ as above.  Then
\[
\alphaJQ(m,n)\le\alphaQ(m,n)+1,
\]
\[
\alphaJS(m,n)\le\alphaQ(m,n)+2,
\]
with the sharper bound $\alphaJS(m,n)\le\alphaQ(m,n)+1$ when $1+m/n$ is integral.  Finally,
\[
f(m,n,L(n))\le(\alphaQ(m,n)+3)m+4n,
\]
and hence
\[
f(m,n,L(n))=O(m\alphaQ(m,n)+n).
\]
\end{corollary}


\appendix

\section{Source-anchor checklist}

This appendix records the source anchors used in the proof.  Page references are to the author PDF / published pagination used in the source-fidelity packet; they are included to make the proof-exposition dependencies auditable.

\begin{itemize}[leftmargin=*]
\item \textbf{Seidel--Sharir, page 7.}  This page contains the rank bound used to justify the logarithmic rank cutoff, the definitions of $g^*$ and $g^\diamond$, and the statement of Lemma~5, the top-down shifting lemma.
\item \textbf{Seidel--Sharir, page 8.}  This page completes the proof of Lemma~5, which is the source-correct shifting lemma with the factor $2n$.
\item \textbf{Seidel--Sharir, page 9.}  This page contains Corollary~6, including the $J$ hierarchy, the recurrence $f(m,n,r)\le km+2nJ_k(r)$, and the source inverse threshold $\alpha_S$ defined through $J_k(\ceil{\log_2 n})\le 1+m/n$.
\item \textbf{Dagstuhl Seminar 25191, Section~5.3.}  This section records the public comparison prompt motivating the direct relationship between the Seidel--Sharir top-down hierarchy and Ackermann-style inverse bounds.
\end{itemize}

The paper uses these anchors only for source setup.  The comparison theorem $R_{z+1}(Q)\ge A(z,4Q)$ is proved internally from the $J$ hierarchy and the packet Ackermann normalization.

\bibliographystyle{alpha}
\bibliography{refs}

\end{document}

